% preamble.tex
% 导言区：宏包、页面设置、代码高亮、超链接等
% 适用于设计赛道文档（docs/design_doc/）

% ==================== 页面设置 ====================
\usepackage{geometry}
\geometry{a4paper, top=2.5cm, bottom=2.5cm, left=3cm, right=2.5cm}

% ==================== 中文支持 ====================
\usepackage[UTF8]{ctex}
% 如果在 Linux 上编译，可能需要指定 fontset:
% \usepackage[UTF8, fontset=ubuntu]{ctex}   % Ubuntu
% \usepackage[UTF8, fontset=founder]{ctex}  % 方正字体

% ==================== 图表和浮动体 ====================
\usepackage{graphicx}
\usepackage{float}
\usepackage{caption}
\usepackage{subcaption}
\usepackage{bm}          % 粗体数学符号
\usepackage{textcomp}    % 特殊符号（如 \textdegree）
\usepackage{upgreek}     % 直立希腊字母

% ==================== 表格增强 ====================
\usepackage{booktabs}    % 三线表
\usepackage{longtable}   % 跨页长表格
\usepackage{array}       % 增强列类型
\usepackage{seqsplit}    % 长序列换行
\usepackage{multirow}    % 合并行
\usepackage{makecell}    % 表头换行
\usepackage{colortbl}    % 表格颜色

% ==================== 数学公式 ====================
\usepackage{amsmath, amssymb, amsfonts}
\usepackage{mathtools}   % amsmath 扩展
\usepackage{amsthm}      % 定理环境

% ==================== 代码高亮 ====================
\usepackage{listings}
\usepackage{xcolor}

% 自定义 listings 样式（移除复杂字符定义）
\lstdefinestyle{codestyle}{
    basicstyle=\ttfamily\small,
    keywordstyle=\color{blue}\bfseries,
    commentstyle=\color{green!40!black},
    stringstyle=\color{red},
    numberstyle=\tiny\color{darkgray},
    numbers=left,
    stepnumber=1,
    breaklines=true,
    breakindent=0pt,
    breakatwhitespace=true,
    showspaces=false,
    showtabs=false,
    tabsize=4,
    frame=single,
    framerule=0.5pt,
    backgroundcolor=\color{white},
    captionpos=b,
    xleftmargin=2mm,
    xrightmargin=2mm,
    belowskip=1.2\baselineskip,
    aboveskip=1.2\baselineskip,
}
\lstset{style=codestyle}

% 为 listings 添加 Rust 语言支持（简化版，移除易出错的 morestring 和 literate）
\lstdefinelanguage{Rust}{
  keywords={
    abstract, as, async, await, break, const, continue, crate, do, dyn, else,
    enum, extern, false, fn, for, if, impl, in, let, loop, match, mod, move,
    mut, pub, ref, return, self, static, struct, super, trait, true, type,
    unsafe, use, where, while, yield,
  },
  keywordstyle=\color{blue}\bfseries,
  ndkeywords={
    Box, Option, Result, String, Vec, u8, u16, u32, u64, u128, usize,
    i8, i16, i32, i64, i128, isize, f32, f64, char, bool,
    &, mut, self, Self, Sync, Send, Arc, Mutex, RwLock,
    Pin, Future, Into, From, TryFrom, TryInto, Default, Clone, Copy, Debug,
    Eq, PartialEq, Ord, PartialOrd, Hash, Deref, DerefMut, Drop,
  },
  ndkeywordstyle=\color{purple}\bfseries,
  identifierstyle=\color{black},
  sensitive=true,
  comment=[l]{//},
  morecomment=[s]{/*}{*/},
  commentstyle=\color{green!40!black}\itshape,
  stringstyle=\color{red},
  morestring=[b]",
  morestring=[b]',
  morestring=[b]",
}

% 对 Rust 代码的额外设置（移除 literate 以避免参数错误）
\lstset{language=Rust}

% ==================== 超链接 ====================
\usepackage{hyperref}
\hypersetup{
    colorlinks=true,
    linkcolor=blue,
    citecolor=blue,
    urlcolor=blue,
    pdfstartview=FitH,
    pdftitle={设计赛道文档},
    pdfauthor={Go\_Go\_Go 参赛队},
    pdfsubject={全国大学生计算机系统能力大赛},
    pdfkeywords={dm-crypt, dm-verity, Rust, Asterinas, 操作系统设计赛}
}

% ==================== 页眉页脚 ====================
\usepackage{fancyhdr}
\pagestyle{fancy}
\fancyhf{}
\fancyhead[L]{\leftmark}
\fancyhead[R]{\thepage}
\renewcommand{\headrulewidth}{0.4pt}
\setlength{\headheight}{15pt}

% 章节页使用 plain 样式（无页眉，有页码）
\fancypagestyle{plain}{
    \fancyhf{}
    \fancyhead[R]{\thepage}
    \renewcommand{\headrulewidth}{0pt}
}

% ==================== 目录格式 ====================
\usepackage{tocloft}
\renewcommand{\cftchapfont}{\bfseries\large}
\renewcommand{\cftchapleader}{\cftdotfill{\cftdotsep}}
\renewcommand{\cftsecfont}{\normalsize}
\renewcommand{\cftsecpagefont}{\normalsize}
\setlength{\cftbeforechapskip}{0.5em}

% ==================== 列表样式 ====================
\usepackage{enumitem}
\setlist{nosep, leftmargin=*}

% ==================== 其他有用宏包 ====================
\usepackage{etoolbox}       % 补丁命令
\usepackage{calc}           % 长度计算
\usepackage{afterpage}      % 在页面后执行命令

% ==================== 图表路径 ====================
% 指向技术文档的 figures 目录，实现图片复用
\graphicspath{{../document/figures/}}

% ==================== 参考文献（可选） ====================
% 如需引用参考文献，取消注释以下行并创建 references.bib
% \usepackage{natbib}
% \bibliographystyle{plain}

% ==================== 章节标题格式 ====================
\usepackage{titlesec}
\titleformat{\chapter}{\Huge\bfseries\centering}{第\thechapter 章}{1em}{}
\titleformat{\section}{\Large\bfseries}{\thesection}{1em}{}
\titleformat{\subsection}{\large\bfseries}{\thesubsection}{1em}{}

% ==================== 附录格式 ====================
% 附录章节编号使用字母
\usepackage{appendix}
\renewcommand{\appendixname}{附录}
\renewcommand{\appendixtocname}{附录}
\renewcommand{\appendixpagename}{附录}

% ==================== 脚注格式 ====================
\renewcommand{\footnotesize}{\small}
\setlength{\footnotesep}{0.5\baselineskip}

% ==================== 段落缩进 ====================
\usepackage{indentfirst}
\setlength{\parindent}{2em}

% ==================== 颜色定义 ====================
\definecolor{codegreen}{rgb}{0,0.6,0}
\definecolor{codegray}{rgb}{0.5,0.5,0.5}
\definecolor{codepurple}{rgb}{0.58,0,0.82}
\definecolor{backcolour}{rgb}{0.95,0.95,0.92}

% ==================== 兼容性设置 ====================
% 防止 listings 与 babel 冲突
\makeatletter
\def\lst@DefSaveDef{%
  \edef\lst@temp{\global\let\csname\@lst @\lst@var\endcsname\@lst@arg}%
  \expandafter\lst@temp
}
\makeatother